\begin{textsample}{POS Dim 6 – ap – Score 15.00 – ap/ap\_em\_46.txt}  \label{ex:f6_pos_101}
Eixos \textbf{Estruturantes} EPT A organização curricular dos cursos técnicos é pautada, em primeiro lugar, pela definição do \textbf{perfil} profissional de conclusão, tornando a concepção curricular um meio pedagógico indispensável na organização dos currículos de educação profissional.
As DNCEM propõem que os itinerários \textbf{formativos} trabalhem unidades curriculares e atividades integradoras, organizadas por meio da preparação básica para o trabalho e de quatro eixos \textbf{estruturantes}, inerentes à ocupação/curso escolhido.
Em conformidade com os eixos \textbf{estruturantes}, nos cursos técnicos foi incluído o componente curricular denominado Projeto Interdisciplinar, que integra conhecimentos voltados para rotinas do eixo tecnológico, o mundo do trabalho e a prática da ocupação estudada. - Investigação - Processos Criativos - \textbf{Mediação} e Intervenção Sociocultural - \textbf{Empreendedorismo} São Eixos \textbf{Estruturantes} : A ) Investigação Científica Integra o componente curricular Projeto Interdisciplinar I, encaminhando o aluno para a realização de procedimentos de investigação voltados à compreensão e enfrentamento de situações, com proposição de intervenções que considerem o desenvolvimento local e a melhoria da qualidade de vida da comunidade.
Define -se o \textbf{problema} de pesquisa e os objetivos, através da estrutura de trabalhos científicos.
Eixo \textbf{Estruturante} : Investigação Científica Competências : - Sistematizar o conhecimento teórico e prático sobre o itinerário \textbf{formativo}, face às exigências \textbf{qualitativas} das condições do \textbf{perfil} do egresso requerido pela sociedade. - Analisar as formas possíveis de investigação científica. - Conhecer e aplicar normas de referenciação de textos científicos ; ler e resumir literatura científica pertinente ao tema escolhido. - Esboçar a estrutura de um projeto, fazendo levantamento de informação e elaboração de resumos e esquemas. - Realizar procedimentos de pesquisa e investigação voltados à compreensão e enfrentamento de situações reais, com proposição de intervenções que considerem o desenvolvimento produtivo local e a melhoria da qualidade de vida da comunidade. - Conhecer a importância e a estrutura de um projeto para realização de atividades interdisciplinares. - Identificar e definir temas para criação do Projeto Interdisciplinar. - Definir o \textbf{problema} da pesquisa, após levantamento das condições de funcionamento do local pesquisado. - Utilizar informações, conhecimentos e ideias \textbf{resultantes} de investigações científicas para criar ou propor soluções para \textbf{problemas} diversos. - Reconhecer os preceitos científicos relacionados aos temas inerentes ao mundo do trabalho. - Utilizar dados, fatos e evidências para \textbf{respaldar} conclusões, opiniões e \textbf{argumentos}, por meio de afirmações claras, \textbf{ordenadas}, coerentes e \textbf{compreensíveis}. - Elaborar hipóteses e perguntas avaliativas sobre o cotidiano do mundo do trabalho e a situação-problema identificada.
B ) Processos Criativos Integra o componente curricular Projeto Interdisciplinar II.
Procura -se expandir a capacidade dos estudantes de idealizar e realizar projetos criativos, a partir dos conhecimentos adquiridos no decorrer do curso.
Eixo \textbf{Estruturante} : Processos Criativos Competências : - Definir e confeccionar os materiais de apoio necessários para execução do projeto. - Discorrer sobre o projeto a ser desenvolvido, destacando os pontos que exigiriam maior atenção. - Demonstrar técnicas de apresentação de projetos. - Registrar ideias, compartilhando experiências pessoais e sugerindo iniciativas. - Utilizar espaços adequados para ideias criativas pertinentes ao estúdio. - Identificar e/ou aprofundar um tema ou \textbf{problema} que orientará a posterior elaboração, apresentação e difusão de uma ação, produto, protótipo, modelo ou solução. - Desenvolver novas abordagens e estratégias para o enfrentamento de situações reais, de forma ética, criativa e inovadora. - Apresentar recursos criativos para resolver \textbf{problemas} reais relacionados a temas e processos de natureza histórica, social, econômica, filosófica, política e/ou cultural, em âmbito local, regional, nacional e/ou global. - Propor soluções éticas, estéticas, criativas e inovadoras para \textbf{problemas} reais relacionados a temas e processos de natureza histórica, social, econômica, filosófica, política e/ou cultural, em âmbito local, regional, nacional e/ou global. - Compreender o impacto e as consequências do uso de novas \textbf{tecnologias} no cotidiano da vida e do trabalho. - Correlacionar formas de expressão das individualidades e singularidades socioemocionais por meio dos saberes propedêuticos e de formação integrada para o mundo do trabalho. - Correlacionar a demanda do \textbf{perfil} profissional com as oportunidades de utilização de recursos criativos inovadores. - Mobilizar intencionalmente recursos criativos relacionados ao conhecimento de projeto, de modo a comunicar com precisão suas ações, reflexões, constatações, interpretações e \textbf{argumentos}.
C ) \textbf{Mediação} e Intervenção Sociocultural No componente curricular Projeto Interdisciplinar III, os estudantes são envolvidos em campos de atuação da vida pública, por meio de projeto que os levem a promover transformações positivas na comunidade, ampliando habilidades relacionadas à convivência e atuação sociocultural ; utilizando esses conhecimentos e habilidades para mediar conflitos, promover entendimentos e propor soluções para questões e \textbf{problemas} socioculturais e \textbf{ambientais} identificados nas comunidades.
Eixo \textbf{Estruturante} : \textbf{Mediação} e Intervenção Sociocultural Competências : - Reconhecer as competências, habilidades e atitudes requeridas para o exercício profissional. - Selecionar intencionalmente conhecimentos e recursos para propor ações individuais e/ou coletivas de \textbf{mediação} e intervenção sobre questões adversas. - Mobilizar ações interventivas de natureza sociocultural e \textbf{ambiental}, em âmbito local, regional e global, baseadas no respeito às diferenças, na escuta, na empatia. - Organizar, promover, coordenar e facilitar o diálogo entre grupos de pessoas e com a comunidade, desenvolvendo técnicas e procedimentos metodológicos adequados. - Aplicar o projeto desenvolvido utilizando conhecimentos das disciplinas e dos Projetos Interdisciplinares I e II. - Utilizar estratégias e métodos para mediar conflitos em equipes heterogêneas, nos diferentes níveis hierárquicos e em diversas e complexas situações, a fim de melhorar o desempenho no local de trabalho. - Promover transformações positivas envolvendo a comunidade, em campos de atuação da vida pública, por meio de iniciativas de mobilização e intervenção sociocultural e \textbf{ambiental}. - Propor medidas que levem a um envolvimento dos indivíduos na vida pública por meio de ações de \textbf{mediação} de conflitos e intervenção sociocultural e \textbf{ambiental}, visando a construção de um espaço de convivência mais harmônico e respeitoso.
D ) \textbf{Empreendedorismo} Com foco no desenvolvimento de processos e produtos com o uso de \textbf{tecnologias} variadas.
O aluno visualiza as características do comportamento empreendedor e sua importância para o desenvolvimento pessoal e profissional, aplica modelos mentais e técnicas de desenvolvimento do \textbf{perfil} empreendedor.
Identifica oportunidades para empreender do mercado, articula competências gerais do curso para construção do Projeto.
Eixo \textbf{Estruturante} : \textbf{Empreendedorismo} Competências : - Analisar o ambiente e avaliar a viabilidade e manutenção de empreendimentos. - Utilizar as características e habilidades de liderança, objetivando o sucesso do projeto. - Reconhecer métodos e técnicas de planejamento estratégico. - Formular propostas concretas, articuladas com o projeto, em âmbito local e/ou global, como forma de intervenção crítica no local estudado. - Propor soluções factíveis para desafios socioculturais e \textbf{ambientais} que afligem a comunidade, com perspectiva de mobilizações locais e/ou regionais. - Desenvolver estratégias éticas, sustentáveis e culturalmente produtivas de concretização do projeto idealizado. - Reconhecer o modelo de gestão da organização, baseando -se na realidade do sistema educacional e considerando a visão sistêmica do empreendimento. - Identificar os instrumentais envolvidos no desenvolvimento dos processos de formulação e gerenciamento do projeto com foco no objetivo pretendido. - Mobilizar ações interventivas de natureza sociocultural/econômica no âmbito local, baseadas no \textbf{empreendedorismo} com responsabilidade socioambiental.
Os eixos \textbf{estruturantes} são complementares e é importante que os itinerários \textbf{formativos} incorporem e integrem todos eles, a fim de garantir que os estudantes experimentem diferentes situações de aprendizagem e desenvolvam um conjunto diversificado de habilidades relevantes para sua formação integral.
A EPT na proposta do Novo Ensino Médio com os Itinerários \textbf{Formativos} tem o poder de contribuir com o aumento da capacidade de (re)inserção social, laboral e política dos seus formandos ; com a extensão de ofertas que contribuam à formação integral dos coletivos que procuram a escola pública de EPT para que esses sujeitos possam atuar, de forma competente e ética, como agentes de mudanças orientadas à satisfação das necessidades coletivas, notadamente as das classes trabalhadoras ( MOURA, 2000 ; FREIRE, 1986 ).
Assim, os estudantes, no decorrer de seu Ensino Médio, deverão realizar pelo menos um itinerário \textbf{formativo} completo, passando necessariamente por todos os eixos.
Cabe salientar que é preciso delinear um desenho curricular em que toda a organização do curso permita o aproveitamento contínuo de estudos, de forma flexível e convergente, articulando as competências e habilidades com : as cargas horárias, as estratégias metodológicas, os critérios de avaliação e os recursos didáticos, equipamentos e materiais.

% matched lemmas: ambiental, argumento, compreensível, empreendedorismo, estruturante, formativo, mediação, ordenado, perfil, problema, qualitativo, respaldar, resultante, tecnologia
\end{textsample}
